% ---------------------------------------------------------------------------
% Appendix: Simulation Procedures
% Source: writeup/contents/appendix.tex:100-225 (reuse, with fixes)
% Fixes applied: (i) appendix.tex:191 "10 independent rounds" -> single-shot with
%   per-configuration repetitions (config-verified: configs/clock/*.yaml, rounds: 1);
% (ii) clock tick stated as \$0.5 per the baseline configs (an earlier draft said
%   "\$1 increments"); (iii) no duplicated figure labels -- prompt-template figures
%   here are unlabeled (the full prompt library, with unique labels, is in app:prompts).
% ---------------------------------------------------------------------------
\section{Simulation Procedures}\label{app:procedures}

We give the prompt-and-parse pipeline used to simulate the auction mechanisms with LLM agents. Each simulated bidder is queried independently. In every auction instance, the agent (i) receives the scenario and mechanism description, (ii) is asked to articulate a short bidding plan, and (iii) outputs a bid/decision in a structured format that can be parsed deterministically. This appendix describes the process for each mechanism class; the complete rule-description prompts for every mechanism and treatment are collected in Appendix~\ref{app:prompts}.

\subsection{Simulation process in sealed-bid auctions}\label{app:sealed-process}

\paragraph{Prompt structure.}
We attach the canonical prompt template in the following. Each bidder is assigned a role (e.g., ``Bidder Andy'') and informed of the other bidders' identities. The prompt then provides (a) a mechanism-specific \emph{rule explanation} (e.g., first-price vs.\ second-price vs.\ third-price), and (b) universal \emph{instructions} that define the objective and response format. We require the model to respond using two exact tags: \verb|<PLAN>| (a brief natural-language rationale) and \verb|<ACTION>| (the bid as a number), in the chain-of-thought style of \citet{wei2022chain}:

\begin{figure}[H]
\centering
\begin{minipage}{0.95\textwidth}
\begin{lstlisting}
    You are Bidder Andy.
    You are bidding with {OTHER_BIDDERS}.

    {RULE EXPLANATION} + {INSTRUCTIONS}

    Your response must use these EXACT tags below:
    ```
    <PLAN>[ Reason out a bidding strategy and plan for this round. {% if current_round > 1 %}Learn from your previous rounds. Remember: the values will be different in future rounds - don't assume patterns will repeat. {% else %}Consider the history above if available. {% endif %}Be detailed and precise. LIMIT your plan to 100 words. ] </PLAN>
    <ACTION> [Your bid amount as a number, e.g., 1, 44, 50] </ACTION>
    ```
\end{lstlisting}
\end{minipage}
\end{figure}

\paragraph{Universal objective instruction.}
Across all sealed-bid simulations, we append a universal instruction block that specifies the optimization objective (maximize profit) and discourages spurious cross-round pattern extrapolation \citep{fish2024algorithmic}. We report the exact text in the following. Note that the template's references to round history are inert in the single-shot baseline (no history block is present in the prompt); they bind only in the repeated-play ablation of Appendix~\ref{app:learning}.

\begin{figure}[H]
\centering
\begin{minipage}{0.95\textwidth}
\begin{lstlisting}
    Your TOP PRIORITY is to place bids which maximize the user's profit in the long run.
    Learn from the history of previous rounds in order to maximize your total profit.
    Don't forget the values are redrawn independently each round.
\end{lstlisting}
\end{minipage}
\end{figure}

\paragraph{History and feedback (multi-round ablation only).}
Our main experiments restrict attention to single-round auctions ($T=1$). When we run the multi-round ablation in Appendix~\ref{app:learning}, we include a standardized history transcript in the prompt. The history contains the agent's past realized values, past bids, whether it won, the transaction price, and realized profit, as well as the list of all bids in each prior round. This transcript is appended above the response-tag requirements, so that the agent can condition on past outcomes without ambiguity.

\subsection{Simulation process in ascending clock auctions}\label{app:clock-process}

In ascending clock auctions, each auction consists of multiple clock cycles, during which every bidder is asked whether they want to stay or drop out at the current clock price.
At the start of the auction, bidders are reminded of the auction rules, similar to those outlined in Appendix~\ref{app:sealed-process}.
The auction starts with an initial price of 0, which increases in fixed increments of \$0.5 per clock tick until only one bidder remains or two bidders drop out simultaneously, in which case the winner is chosen randomly.\footnote{The \$0.5 tick is taken from the baseline experiment configuration files (\texttt{configs/clock/}); an earlier draft stated \$1 increments, which does not match any configuration. Some model-ablation configurations used a \$0.1 tick instead of \$0.5. \TODO{Harmonize all clock cells to the canonical \$0.5 tick of Table~\ref{tab:calibration-master} in the re-runs (plan E2).}}
The detailed prompt is listed as follows, with variables enclosed in brackets (the value shown, 26, is a sample instantiation):

\begin{figure}[H]
\centering
\begin{minipage}{0.95\textwidth}
\begin{lstlisting}
    You are Bidder Andy. You are bidding with {OTHER_BIDDERS}.

    {RULE EXPLANATION} + {INSTRUCTIONS}

    Your value towards to the prize is 26 in this round.
    The current price in this clock cycle is {current_price}.
    The price for next clock cycle is {current_price + increment}.

    The previous bidding history is: {transcript}.

    Do you want to stay in the bidding?
    If you choose yes, you can keep bidding for next clock. If you choose No, you will exit and have no chance to re-enter the bidding. Your response must use these EXACT tags below. You must output the ACTION.
    ```
    <PLAN>
    [Write your plans for bidding strategies. Be detailed and precise but keep things succinct and don't repeat yourself. LIMIT your plan to 50 words. ] </PLAN>
    <ACTION> Yes or No </ACTION>
    ```
\end{lstlisting}
\end{minipage}
\end{figure}

In AC (the open clock), bidders who have not dropped out are shown the previous bidding history at each subsequent clock cycle. The history includes the price and how many bidders had dropped out,
and the transcript variable takes the following format:
\textit{The previous biddings are:
['In clock round 1, the price was 1,
no players dropped out'].}

In AC-B (the closed/blind clock), we do not show the previous biddings; in each clock cycle, we directly query the LLM agents' decision,
so the transcript variable is None.

A clock auction ends when only one bidder remains, who is declared the winner, or when two bidders drop out simultaneously, in which case the winner is chosen randomly. Each clock auction is a \emph{single-shot} instance with affiliated values (\texttt{rounds:\,1} in the configuration files): agents receive no feedback across instances, and we run $K$ independent repetitions per configuration with freshly drawn values ($K=40$ in the baseline clock cells; config-verified). \TODO{Top up the clock cells to the pinned $K \geq 50$ of the harmonized design in the re-runs (plan \S 4, E2).}

\subsection{Simulation process in eBay auctions}\label{app:ebay-process}

For eBay auctions, we discretize the continuous bidding period into 10 periods. Each period, LLM agents decide whether to increase their bid or hold their current bid. The bidding process follows a structured format, where bidders act in a predefined sequence each period (e.g., Charles, then Alice, then Betty). Agents are informed of past price changes through a transcript, such as ``On day 1, the price changed to 1. On day 2, the price changed to 3.''

The prompt provides key details, including the bidder's private value, the total number of bidding days, the current day, the bidding order, previous bids, previous proxy bids, and the current price. The decision-making process is again guided by a structured asking prompt that enforces bidding rules. The treatment-level analysis of this environment appears in Appendix~\ref{app:ebay}.

\begin{figure}[H]
\centering
\begin{minipage}{0.95\textwidth}
\begin{lstlisting}
    You are Bidder Andy.
    You are bidding with {OTHER_BIDDERS}.

    {RULE EXPLANATION} + {INSTRUCTIONS}

    Your private value for this item is ${private_value}. This is the maximum amount you are willing to pay. Keep this value private. There are in total {total_periods} days  of bidding and this is day {current_period}. {ordering}.

    Your previous bids are {previous_bid}. If you have already placed bids, you can only increase your bid or hold your current bid. The previous proxy bids are: {transcript}.
    The current price is {current_price}.
    Your response must use these EXACT tags below. Don't miss the amount.
    ```
    <PLAN>[Write your plans for what bidding strategies to test next. Be detailed and precise but keep things succinct and don't repeat yourself. LIMIT your plan to 50 words. ]</PLAN>
    <CHECK>your last bid is {{last_bid_amount}}, you cannot bid lower that this value </CHECK>
    <ACTION> BID or HOLD  </ACTION>
    <AMOUNT> if BID, enter a number here, e.g. 1, 44. If HOLD, enter 0 </AMOUNT>
    ```
\end{lstlisting}
\end{minipage}
\end{figure}
